\documentclass[12pt,a4paper]{article}
\usepackage[utf8]{inputenc}
\usepackage[T1]{fontenc}
\usepackage[french]{babel}
\usepackage{geometry}
\usepackage{hyperref}
\usepackage{listings}
\usepackage{xcolor}
\usepackage{fontspec}
\usepackage{fontawesome}
\usepackage{fontawesome5}
\setsansfont{Noto Sans}[Scale=MatchLowercase]

\geometry{margin=2cm}

\lstset{
    language=HTML,
    basicstyle=\ttfamily\small,
    keywordstyle=\color{blue},
    commentstyle=\color{green},
    stringstyle=\color{red},
    showstringspaces=false,
    frame=single,
    breaklines=true
}

\title{Guide des icônes FontAwesome}
\author{Pour l'activité CSS}
\date{\today}

\begin{document}

\maketitle

\section{Introduction}

Les icônes utilisées dans le site marchand sont issues de la bibliothèque FontAwesome. Cette bibliothèque fournit des milliers d'icônes vectorielles qui peuvent être intégrées facilement dans vos pages HTML.

\section{Installation de FontAwesome}

Pour utiliser FontAwesome dans votre projet, vous devez inclure le fichier CSS de FontAwesome dans votre page HTML :

\begin{lstlisting}
<head>
    <link rel="stylesheet" href="https://cdnjs.cloudflare.com/ajax/libs/font-awesome/6.0.0/css/all.min.css">
</head>
\end{lstlisting}

\section{Utilisation des icônes}

\subsection{Syntaxe de base}

Les icônes FontAwesome s'utilisent avec des classes CSS spécifiques :

\begin{lstlisting}
<i class="fas fa-shopping-cart"></i>  <!-- Icône panier plein -->
<i class="far fa-shopping-cart"></i>  <!-- Icône panier vide -->
<i class="fa fa-cookie"></i>           <!-- Icône cookie -->
\end{lstlisting}

\subsection{Structure des classes}

\begin{itemize}
\item \texttt{fas} : FontAwesome Solid (icônes pleines)
\item \texttt{far} : FontAwesome Regular (icônes vides)
\item \texttt{fab} : FontAwesome Brands (marques/logos)
\item \texttt{fa-} suivi du nom de l'icône
\end{itemize}

\section{Icônes utilisées dans le projet}

Voici les principales icônes utilisées dans le site marchand :

\subsection{Icônes de navigation et commerce}

\begin{itemize}
\item \faShoppingCart{} \texttt{<i class="fas fa-shopping-cart"></i>} : Panier d'achats
\item \faHome{} \texttt{<i class="fas fa-home"></i>} : Page d'accueil
\item \faArchive{} \texttt{<i class="fas fa-archive"></i>} : Produits/Archive
\item \faPencil{} \texttt{<i class="fas fa-pencil"></i>} : Inscription/Édition
\item \faMap{} \texttt{<i class="fas fa-map"></i>} : Plan du site
\end{itemize}

\subsection{Icônes d'information}

\begin{itemize}
\item \faInfo{} \texttt{<i class="fas fa-info"></i>} : Informations
\item \faClockO{} \texttt{<i class="far fa-clock"></i>} : Horloge/Temps
\item \faCookie{} \texttt{<i class="fas fa-cookie"></i>} : Cookies
\end{itemize}

\subsection{Icônes de recherche}

\begin{itemize}
\item \faSearch{} \texttt{<i class="fas fa-search"></i>} : Recherche
\end{itemize}

\subsection{Icônes mobiles}

\begin{itemize}
\item \faMobile{} \texttt{<i class="fas fa-mobile"></i>} : Téléphone mobile
\end{itemize}

\section{Exemples d'utilisation dans le HTML}

\subsection{Dans le titre principal}

\begin{lstlisting}
<h1><i class="fas fa-shopping-cart"></i> Boutique SNT</h1>
\end{lstlisting}

\subsection{Dans les liens de navigation}

\begin{lstlisting}
<nav>
    <a href="index.html"><i class="fas fa-home"></i> Accueil</a>
    <a href="pages/produits.html"><i class="fas fa-archive"></i> Produits</a>
    <a href="pages/panier.html"><i class="fas fa-shopping-cart"></i> Panier</a>
</nav>
\end{lstlisting}

\subsection{Dans les boutons}

\begin{lstlisting}
<button>
    <i class="fas fa-cookie"></i> Voir mes cookies
</button>
\end{lstlisting}

\subsection{Dans les titres de section}

\begin{lstlisting}
<h2><i class="far fa-clock"></i> Vous avez récemment consulté</h2>
\end{lstlisting}

\section{Styles CSS pour les icônes}

Vous pouvez styliser les icônes avec du CSS :

\begin{lstlisting}
.fa-shopping-cart {
    color: #1abc9c;        /* Couleur personnalisée */
    font-size: 1.2em;      /* Taille plus grande */
    margin-right: 5px;    /* Espace à droite */
}

.fa-cookie:hover {
    color: #e74c3c;        /* Changement de couleur au survol */
    transform: scale(1.1); /* Agrandissement au survol */
}
\end{lstlisting}

\section{Trouver d'autres icônes}

Pour découvrir toutes les icônes disponibles :

\begin{enumerate}
\item Rendez-vous sur le site officiel : \url{https://fontawesome.com/icons}
\item Utilisez la barre de recherche pour trouver l'icône souhaitée
\item Copiez le nom de la classe (ex: \texttt{fa-shopping-cart})
\item Choisissez le style approprié (solid, regular, brands)
\end{enumerate}

\section{Conseils pratiques}

\begin{itemize}
\item Les icônes s'adaptent automatiquement à la taille du texte parent
\item Vous pouvez changer leur couleur avec la propriété CSS \texttt{color}
\item Utilisez \texttt{font-size} pour changer leur taille
\item Les icônes sont vectorielles : elles restent nettes à toutes les tailles
\item Pensez à l'accessibilité : ajoutez un attribut \texttt{aria-label} si nécessaire
\end{itemize}

\section{Exemple complet}

Voici un exemple complet d'utilisation des icônes dans une page :

\begin{lstlisting}
<!DOCTYPE html>
<html lang="fr">
<head>
    <meta charset="UTF-8">
    <title>Ma Boutique</title>
    <link rel="stylesheet" href="https://cdnjs.cloudflare.com/ajax/libs/font-awesome/6.0.0/css/all.min.css">
    <link rel="stylesheet" href="styles.css">
</head>
<body>
    <header>
        <h1><i class="fas fa-shopping-cart"></i> Ma Boutique</h1>
        <nav>
            <a href="#"><i class="fas fa-home"></i> Accueil</a>
            <a href="#"><i class="fas fa-archive"></i> Produits</a>
            <a href="#"><i class="fas fa-shopping-cart"></i> Panier</a>
        </nav>
    </header>

    <main>
        <section>
            <h2><i class="far fa-clock"></i> Produits récents</h2>
            <button><i class="fas fa-cookie"></i> Voir cookies</button>
        </section>
    </main>
</body>
</html>
\end{lstlisting}

Bonne chance avec votre projet !

\end{document}